\documentclass[11pt,a4paper]{article}

\usepackage[provide=*,spanish]{babel}
\usepackage[utf8]{inputenc}
\usepackage[T1]{fontenc}
\usepackage{lmodern}
\usepackage[margin=2.35cm]{geometry}
\usepackage{amsmath}
\usepackage{booktabs}
\usepackage{longtable}
\usepackage{tabularx}
\usepackage{array}
\usepackage{enumitem}
\usepackage{listings}
\usepackage{xcolor}
\usepackage{microtype}
\usepackage{parskip}
\usepackage{titlesec}
\usepackage{hyperref}
\usepackage{caption}
\usepackage{seqsplit}

\definecolor{azulWCD}{RGB}{17,67,112}
\definecolor{azulClaro}{RGB}{232,241,249}
\definecolor{grisCodigo}{RGB}{247,247,247}
\definecolor{verdeCodigo}{RGB}{25,110,55}
\definecolor{rojoCodigo}{RGB}{145,35,35}
\definecolor{verdeDiff}{RGB}{20,110,50}
\definecolor{rojoDiff}{RGB}{155,25,25}
\definecolor{azulDiff}{RGB}{40,80,150}
\definecolor{tituloColor}{RGB}{0,55,110}
\definecolor{reglaColor}{RGB}{0,90,160}

\hypersetup{
  colorlinks=true,
  linkcolor=azulWCD,
  urlcolor=azulWCD,
  citecolor=azulWCD,
  pdftitle={Integracion completa de las mejoras WCD/NaCl en MEIGA},
  pdfauthor={Jaime A. Betancourt}
}

\lstdefinestyle{cppstyle}{
  language=C++,
  backgroundcolor=\color{grisCodigo},
  basicstyle=\ttfamily\footnotesize,
  keywordstyle=\color{azulWCD}\bfseries,
  commentstyle=\color{verdeCodigo}\itshape,
  stringstyle=\color{rojoCodigo},
  numbers=left,
  numberstyle=\tiny\color{gray},
  numbersep=8pt,
  breaklines=true,
  breakatwhitespace=false,
  frame=single,
  rulecolor=\color{gray!40},
  tabsize=2,
  showstringspaces=false,
  xleftmargin=1em,
  framexleftmargin=1em,
  columns=fullflexible,
  keepspaces=true
}

\lstdefinestyle{diffstyle}{
  backgroundcolor=\color{grisCodigo},
  basicstyle=\ttfamily\scriptsize,
  breaklines=true,
  breakatwhitespace=false,
  frame=single,
  rulecolor=\color{gray!40},
  tabsize=2,
  showstringspaces=false,
  xleftmargin=1em,
  framexleftmargin=1em,
  columns=fullflexible,
  keepspaces=true,
  literate=
    {á}{{\'a}}1 {é}{{\'e}}1 {í}{{\'i}}1 {ó}{{\'o}}1 {ú}{{\'u}}1
    {ñ}{{\~n}}1 {Á}{{\'A}}1 {É}{{\'E}}1 {Í}{{\'I}}1 {Ó}{{\'O}}1
    {Ú}{{\'U}}1 {Ñ}{{\~N}}1 {®}{{\textregistered}}1
}
\lstset{style=cppstyle}

\titleformat{\section}{\normalfont\Large\bfseries\color{azulWCD}}
  {\thesection.}{0.6em}{}
\titleformat{\subsection}{\normalfont\large\bfseries}
  {\thesubsection.}{0.5em}{}
\setlist{itemsep=0.25em,topsep=0.4em}
\renewcommand{\arraystretch}{1.18}
\setlength{\emergencystretch}{2em}

\newcommand{\codigo}[1]{\nolinkurl{#1}}
\newcommand{\archivo}[1]{\path{#1}}
\newcommand{\archnom}[1]{\ttfamily\seqsplit{#1}}

\begin{document}
\pagestyle{empty}

%------------------------------------------------
% PORTADA
%------------------------------------------------

\begin{titlepage}
\begin{center}

\vspace*{0.5cm}
\textcolor{reglaColor}{\rule{\linewidth}{1.2pt}}

\vspace{1.3cm}
{\Large\scshape Universidad Industrial de Santander}\\[0.2cm]
{\large Escuela de Física --- Facultad de Ciencias}\\[0.2cm]
{\large Doctorado en Física}

\vspace{2.2cm}
{\Huge\bfseries\color{tituloColor} Integración completa de las\\[0.2cm]
mejoras del WCD con soluciones\\[0.2cm] de NaCl en MEIGA}

\vspace{1.8cm}
{\large\textbf{Informe técnico}}

\vspace{0.6cm}
\begin{minipage}{0.82\linewidth}
\centering
\textit{``Variante \texttt{MEIGA-WCD-SalphaBeta-2026}''}
\end{minipage}

\vfill

{\large\textbf{Autor}}\\[0.15cm]
Jaime A. Betancourt\\[0.1cm]
Doctorado en Física, Universidad Industrial de Santander

\vspace{1.2cm}
\textcolor{reglaColor}{\rule{0.5\linewidth}{0.6pt}}

\vspace{1cm}
13 de septiembre de 2026

\vspace{1cm}
\textcolor{reglaColor}{\rule{\linewidth}{1.2pt}}

\end{center}
\end{titlepage}

\pagestyle{plain}

\begin{abstract}
\noindent
Este informe documenta la integración, en una única base de código
consistente, de las cinco familias de mejoras desarrolladas de forma
incremental para el detector Cherenkov de agua (WCD) del framework MEIGA:
(1) el registro paso a paso de \codigo{G4WCDSteppingAction}; (2) las
densidades y la composición de las soluciones de NaCl al 2.5, 5.0 y 10.0\,\%;
(3) la corrección óptica del recubrimiento de Tyvek (propiedad
\codigo{RINDEX} en \codigo{linerOpticalPT}) junto con las tablas ópticas
diferenciadas por concentración; (4) una propuesta alternativa de
propiedades ópticas basada en interpolación y dispersión de Rayleigh; y
(5) las correcciones de la respuesta del PMT (\codigo{G4MPMTAction},
\codigo{OptDevice}, material efectivo del fotocátodo). El resultado es un
único árbol de código fuente, derivado de la variante
\codigo{MEIGA-WCD-SalphaBeta-2026}, que conserva el tratamiento térmico
S($\alpha,\beta$) del hidrógeno del agua y añade, de forma consistente,
todas las correcciones ópticas y de respuesta del PMT revisadas en los
informes técnicos previos. Se documentan también las inconsistencias
encontradas en los archivos de entrada (una rama óptica alternativa
incompatible con la finalmente adoptada, un archivo con nombre incorrecto,
y un archivo huérfano con nombre corrupto) y la manera en que se
resolvieron.
\end{abstract}

\tableofcontents

\section{Objeto y alcance del informe}

El propósito de este documento es dejar constancia completa, reproducible
y auditable de los cambios de código fuente realizados sobre la variante
\codigo{MEIGA-WCD-SalphaBeta-2026} del framework MEIGA para incorporar las
cinco familias de mejoras desarrolladas para el estudio del WCD con agua y
soluciones de NaCl. El informe:

\begin{itemize}
  \item identifica el archivo base (\emph{Meiga anterior}) sobre el cual se
    aplicaron los cambios;
  \item describe la metodología usada para reconstruir, a partir de once
    archivos \codigo{*-mejorado}/\codigo{*-corregido} parcialmente
    superpuestos, una única versión consistente de cada archivo fuente;
  \item documenta cada cambio, archivo por archivo, con su motivación
    física o de ingeniería;
  \item reporta explícitamente los hallazgos que requirieron una decisión
    de integración (rama óptica alternativa no incorporada, archivo mal
    etiquetado, archivo huérfano); y
  \item incluye, en el Apéndice~\ref{ap:diffs}, el \emph{diff} unificado
    completo (archivo original de MEIGA contra archivo final integrado)
    de cada uno de los ocho archivos modificados.
\end{itemize}

No se recompiló el simulador dentro de este entorno de trabajo porque no
hay una instalación de Geant4 disponible; la validación realizada aquí es
estática (revisión de consistencia de símbolos, balance de llaves y
paréntesis, y verificación cruzada de todo el árbol \codigo{src/} en busca
de referencias rotas). La Sección~\ref{sec:validacion} detalla el alcance
exacto de esta verificación y lo que queda pendiente en la máquina de
simulación.

\section{Archivo base y archivos de entrada}

\subsection{Base: \emph{Meiga anterior}}

El punto de partida es el árbol de código fuente contenido en
\codigo{MEIGA-WCD-SalphaBeta-2026.zip}, una variante independiente de MEIGA
que ya incorpora el tratamiento térmico S($\alpha,\beta$) del hidrógeno
ligado del agua (elemento \codigo{TS\_H\_of\_Water}, registro de
\codigo{G4ThermalNeutrons} sobre \codigo{QGSP\_BERT\_HP} en
\codigo{G4MPhysicsList.cc}, y las cuatro densidades de referencia a
20\,\textdegree C: 0.99820, 1.01604, 1.03401 y 1.07051\,g/cm\textsuperscript{3}
para 0, 2.5, 5 y 10\,\% de NaCl). Este es el árbol referido como
\emph{Meiga anterior} en el resto del informe.

\subsection{Archivos de mejora aportados}

Se recibieron once archivos fuente \codigo{*-mejorado}/\codigo{*-corregido}
distribuidos en cuatro carpetas de trabajo, además de sus informes
técnicos en \LaTeX. La Tabla~\ref{tab:insumos} resume su procedencia y su
relación de dependencia, reconstruida a partir de las marcas de tiempo del
sistema de archivos y de un análisis diferencial exhaustivo entre cada par
de versiones sucesivas.

\begin{longtable}{@{}p{4.0cm}p{1.9cm}p{7.9cm}@{}}
\toprule
\textbf{Archivo} & \textbf{Grupo} & \textbf{Contenido / relación} \\
\midrule
\endhead
{\small\archnom{G4WCDSteppingAction-mejorado.h/.cc}} & 1 & Reescritura completa e independiente. \\
{\small\archnom{Materials-mejorado.cc}, \archnom{SaltyWCD-mejorado.cc}} & 2 & Primera revisión: densidades a 20\,\textdegree C y las tres soluciones discretas de NaCl. Base de las revisiones 3 y 4. \\
{\small\archnom{Materials-mejorado-optica-Tyveck.cc}, \archnom{SaltyWCD-mejorado-optica-Tyveck.cc}} & 3a & Deriva de la revisión 2: tabla óptica \codigo{RINDEX/ABSLENGTH/MIEHG} propia por concentración. \\
{\small\archnom{Materials-mejorado-optica-Tyvek-RINDEX.cc}, \archnom{SaltyWCD-mejorado-optica-Tyvek-RINDEX.cc}} & 3b & Deriva de 3a: añade \codigo{RINDEX} a \codigo{linerOpticalPT} y cambia la superficie del Tyvek a \codigo{dielectric\_dielectric}/\codigo{groundbackpainted}. \\
{\small\archnom{Materials-mejorado-optica.cc}, \archnom{SaltyWCD-mejorado-optica.cc}} & 4 & Deriva de la revisión 2, en una rama \textbf{independiente} de 3a/3b: modelo de interpolación normalizado al índice en la línea D del sodio, con dispersión de Rayleigh. Cronológicamente anterior a 3a y \textbf{no continuada}. \\
{\small\archnom{Materials-corregido.cc}, \archnom{SaltyWCD-corregido.cc}} & 5 & Deriva de 3b: añade las correcciones de material efectivo del PMT (\codigo{Pyrex}) y la geometría por capas explícitas del WCD. \textbf{Versión final} de \codigo{Materials.cc}/\codigo{SaltyWCD.cc}. \\
{\small\archnom{G4MPMTAction-mejorado.h/.cc}} & 5 & Activación en \codigo{fGeomBoundary}, terminación del fotón, deduplicación de fotoelectrones. \\
{\small\archnom{OptDevice-mejorado.cc}} & 5 & Interpolación de la eficiencia cuántica R5912 (300--650\,nm), remoción de archivos de depuración por fotón. \\
{\small\archnom{Materials.h} (en \archnom{Correcciones-PMT/})} & --- & \textbf{Hallazgo}: el contenido de este archivo es en realidad una copia intermedia de \codigo{SaltyWCD::BuildDetector}, no un encabezado. Ver Sección~\ref{sec:hallazgos}. \\
{\small\archnom{Detector.h/.cc}, \archnom{OptDevice.h}, \archnom{OpticalPhysics.cc/.hh}, \archnom{G4MPMTAction.h/.cc} (sin sufijo, en \archnom{Correcciones-PMT/})} & --- & Copias de referencia, \textbf{idénticas byte a byte} al Meiga anterior. No requieren ninguna acción. \\
\bottomrule
\caption{Procedencia y relación de dependencia de los archivos de mejora recibidos.}
\label{tab:insumos}
\end{longtable}

\section{Metodología de integración}
\label{sec:metodologia}

La integración se realizó en cuatro pasos:

\begin{enumerate}
  \item \textbf{Reconstrucción del árbol de dependencias.} Se calculó el
    \emph{hash} MD5 de cada archivo para detectar duplicados exactos entre
    carpetas (por ejemplo, \codigo{G4WCDSteppingAction-mejorado.cc} aparece
    igual en la carpeta raíz y en \codigo{Correcciones-PMT/}), y se generó
    el \emph{diff} unificado entre cada par de versiones sucesivas de
    \codigo{Materials.cc} y \codigo{SaltyWCD.cc} para determinar qué
    archivo es una extensión directa de cuál otro (Tabla~\ref{tab:insumos}).
  \item \textbf{Selección de la rama a integrar.} Cuando dos archivos
    resultaron ser implementaciones alternativas e incompatibles de la
    misma funcionalidad (grupo 3 frente a grupo 4, ambos sobre
    propiedades ópticas diferenciadas por concentración de NaCl), se
    adoptó la rama que efectivamente continuó hacia la versión corregida
    final (grupo 3 $\to$ grupo 5), y se documentó la rama no adoptada
    (Sección~\ref{sec:grupo4}).
  \item \textbf{Reintegración del tratamiento S($\alpha,\beta$).} Se
    verificó, comparando contra el Meiga anterior, que la cadena de
    archivos de mejora fue desarrollada a partir de una copia de
    \codigo{Materials.cc} \emph{sin} el elemento
    \codigo{TS\_H\_of\_Water}. Se restituyó este elemento y su uso en
    \codigo{Water} y en las tres soluciones salinas, de modo que el
    simulador final conserve el modelo térmico ya validado en el Meiga
    anterior (Sección~\ref{sec:salphabeta}).
  \item \textbf{Verificación de consistencia cruzada.} Se recorrió todo
    \codigo{src/} en busca de referencias a los símbolos modificados
    (\codigo{Materials::SaltyWater}, \codigo{Materials::HDPE},
    \codigo{Materials::Pyrex}, \codigo{Materials::LinerOptSurf2}, nombres
    de materiales usados en búsquedas por cadena, etc.) para confirmar que
    ningún otro modelo de detector (\codigo{WCD.cc}, \codigo{Hodoscope.cc},
    \codigo{Mudulus.cc}, \codigo{Musaic.cc}, \codigo{Scintillator.cc}) ni
    ninguna aplicación quedara con una referencia rota.
\end{enumerate}

\subsection{Hallazgos durante la integración}
\label{sec:hallazgos}

\begin{itemize}
  \item \textbf{Archivo \codigo{Correcciones-PMT/Materials.h} mal
    etiquetado.} Su contenido (424 líneas) no corresponde a un encabezado
    de clase: es, byte a byte salvo una única diferencia menor, una
    versión intermedia de \codigo{SaltyWCD::BuildDetector} (idéntica en
    un 99\,\% a \codigo{SaltyWCD-corregido.cc}). No se usó como fuente de
    cambios para \codigo{Materials.h}; en su lugar, las declaraciones
    necesarias para el nuevo encabezado (\codigo{SaltyWater\_2p5},
    \codigo{SaltyWater\_5}, \codigo{SaltyWater\_10}) se derivaron
    directamente de los símbolos nuevos que \codigo{Materials-corregido.cc}
    efectivamente define y usa.
  \item \textbf{Archivo huérfano con nombre corrupto.} El paquete
    \codigo{MEIGA-WCD-SalphaBeta-2026.zip} contenía, en
    \codigo{src/G4Models/}, un archivo cuyo nombre es el carácter suelto
    «ç» (sin extensión). Su contenido es un borrador antiguo de
    \codigo{SaltyWCD::BuildDetector} anterior incluso a la introducción del
    modelo S($\alpha,\beta$) (usa \codigo{new G4Material("SaltyWater",
    1.1*g/cm3, 2)}, sin relación con las densidades a 20\,\textdegree C
    documentadas en el CHANGELOG). Al no tener extensión \codigo{.cc}, el
    sistema de compilación (\codigo{file(GLOB sources *.cc)} en
    \codigo{CMakeLists.txt}) nunca lo incluyó en la compilación, por lo
    que es inerte pero confuso. Se eliminó del árbol final.
  \item \textbf{Rama óptica alternativa (grupo 4) no continuada.} Ver
    Sección~\ref{sec:grupo4}.
\end{itemize}

\section{Resumen de archivos modificados}

\begin{longtable}{@{}p{6.6cm}p{1.5cm}p{1.5cm}p{3.6cm}@{}}
\toprule
\textbf{Archivo (ruta relativa a \codigo{src/})} & \textbf{Líneas orig.} & \textbf{Líneas final} & \textbf{Grupo(s)} \\
\midrule
\endhead
{\small\archnom{Applications/G4WCDSimulator/G4WCDSteppingAction.h}} & 44 & 33 & 1 \\
{\small\archnom{Applications/G4WCDSimulator/G4WCDSteppingAction.cc}} & 398 & 338 & 1 \\
{\small\archnom{G4Models/Materials.h}} & 115 & 120 & 2, 3, 5 \\
{\small\archnom{G4Models/Materials.cc}} & 654 & 875 & 2, 3, 5 (+ S($\alpha,\beta$)) \\
{\small\archnom{G4Models/SaltyWCD.h}} & 34 & 34 & sin cambios \\
{\small\archnom{G4Models/SaltyWCD.cc}} & 258 & 466 & 2, 3, 5 \\
{\small\archnom{G4Models/G4MPMTAction.h}} & 37 & 38 & 5 \\
{\small\archnom{G4Models/G4MPMTAction.cc}} & 141 & 144 & 5 \\
{\small\archnom{Framework/Detector/OptDevice.h}} & 155 & 155 & sin cambios \\
{\small\archnom{Framework/Detector/OptDevice.cc}} & 363 & 341 & 5 \\
\bottomrule
\caption{Archivos de código fuente afectados por la integración.}
\label{tab:resumen}
\end{longtable}

\section{Descripción detallada de los cambios por grupo}

\subsection{\texorpdfstring{Grupo 1 --- Registro paso a paso: \codigo{G4WCDSteppingAction}}{Grupo 1 --- Registro paso a paso: G4WCDSteppingAction}}

La clase se reescribió por completo manteniendo su interfaz pública
(constructor y \codigo{UserSteppingAction}). Los cambios principales son:

\begin{itemize}
  \item \textbf{Unidades explícitas y normalizadas}: todas las longitudes
    se expresan en cm y las energías en MeV o eV según corresponda,
    eliminando la mezcla de unidades del Meiga anterior.
  \item \textbf{Identificadores de corrida y evento}: cada fila TSV incluye
    \codigo{run\_id} y \codigo{event\_id}, obtenidos de
    \codigo{G4RunManager::GetCurrentRun()/GetCurrentEvent()}.
  \item \textbf{Identificación correcta del núcleo capturador}: en vez de
    leer \codigo{material->GetElement(0)->GetName()} (que en una mezcla
    reporta arbitrariamente el primer elemento declarado, no el que
    realmente participó en la reacción), se usa
    \codigo{G4HadronicProcess::GetTargetNucleus()} para obtener el
    (Z, A) real del núcleo capturador.
  \item \textbf{Separación primarias/secundarias}: las tablas
    \codigo{neutrones-incidentes.tsv}, \codigo{secundarias-captura-\allowbreak
    neutron.tsv}, \codigo{secundarias-inelastica-neutron.tsv},
    \codigo{secundarias-interaccion-gamma.tsv} y
    \codigo{fotones-cherenkov.tsv} separan explícitamente el registro de
    la partícula primaria (paso 1, \codigo{parentID == 0}) del de las
    partículas secundarias generadas en cada interacción.
  \item \textbf{Escritura seguridad en multihilo}: \codigo{AppendTsv} usa
    un \codigo{G4Mutex} (\codigo{G4AutoLock}) para serializar la escritura
    concurrente de varios hilos de Geant4 sobre el mismo archivo TSV, y
    escribe el encabezado una única vez.
  \item \textbf{Formato}: los antiguos archivos \codigo{.txt} de columnas
    separadas por tabulador sin encabezado se reemplazan por tablas
    \codigo{.tsv} con una línea de encabezado explícita
    (\codigo{stepHeader}, \codigo{secondaryHeader}, \codigo{captureHeader}).
\end{itemize}

\subsection{Grupo 2 --- Densidades y composición de las soluciones de NaCl}

Se añaden a \codigo{Materials} tres materiales discretos,
\codigo{SaltyWater\_2p5}, \codigo{SaltyWater\_5} y \codigo{SaltyWater\_10},
construidos por \codigo{CreateSaltyWaterByFinalMassFraction} a partir de la
fracción másica final $w_{\text{NaCl}} = m_{\text{NaCl}} / (m_{\text{H}_2\text{O}} + m_{\text{NaCl}})$,
con las densidades a 20\,\textdegree C ya validadas en el Meiga anterior
(0.99820, 1.01604, 1.03401 y 1.07051\,g/cm\textsuperscript{3}). El
material anterior \codigo{Materials::SaltyWater} (una única mezcla
continua construida dinámicamente en \codigo{SaltyWCD::BuildDetector} para
cualquier fracción de NaCl) se conserva como alias de compatibilidad hacia
\codigo{SaltyWater\_10}, pero \codigo{SaltyWCD::BuildDetector} ahora
selecciona explícitamente uno de los cuatro medios predefinidos
(\codigo{Water}, \codigo{SaltyWater\_2p5/5/10}) mediante la función
\codigo{SelectWCDMedium}, que aborta la ejecución con
\codigo{G4Exception} si se solicita una concentración no tabulada, en
lugar de aproximarla silenciosamente. También se corrigen, en esta misma
revisión, tres propiedades ópticas del Tyvek que estaban físicamente mal
definidas en el Meiga anterior: la longitud de absorción de vol\'umen
pasa de 10\,m (transparente) a 0.1\,mm (opaco, consistente con un
reflector difuso real), el índice de refracción de vol\'umen pasa de 0 a
1.49, y el \codigo{BACKSCATTERCONSTANT} pasa de 0 a 0.02.

\subsection{Grupo 3 --- Corrección óptica del Tyvek y tablas por concentración}

Esta revisión reemplaza la única tabla óptica compartida
\codigo{waterPT1}/\codigo{waterPT2} (originalmente un índice de refracción
constante de dos puntos, 1.33 en todo el espectro) por una tabla completa
e independiente por medio activo (\codigo{CreateWCDOpticalTable}), cada una
con su propio \codigo{RINDEX} espectral de 32 puntos (agua pura y las tres
soluciones), una \codigo{ABSLENGTH} común como línea base, y el modelo Mie
del ejemplo \codigo{OpNovice} de Geant4 conservado explícitamente como
aproximación declarada. Adicionalmente, se corrige
\codigo{linerOpticalPT}: se le añade la propiedad \codigo{RINDEX}
(requerida por \codigo{G4OpBoundaryProcess} para una superficie
\codigo{dielectric\_dielectric}), y la superficie \codigo{LinerOptSurf}
pasa de \codigo{dielectric\_metal}/\codigo{ground} a
\codigo{dielectric\_dielectric}/\codigo{groundbackpainted}, el modelo
correcto para un recubrimiento reflectante sobre un medio dieléctrico
(agua) en vez de un metal.

\subsection{Grupo 4 --- Rama óptica alternativa (no integrada)}
\label{sec:grupo4}

\codigo{Materials-mejorado-optica.cc} y \codigo{SaltyWCD-mejorado-optica.cc}
proponen un modelo distinto para la misma necesidad física (diferenciar
las propiedades ópticas del agua pura y las soluciones de NaCl): en vez de
tablas \codigo{RINDEX} medidas punto a punto, interpolan la curva de
dispersión del agua pura y la desplazan verticalmente hasta que su valor
en la línea D del sodio (589.3\,nm) coincide con un índice objetivo por
concentración; añaden además dispersión de Rayleigh explícita
(\codigo{RAYLEIGH}) con una longitud de referencia de 100\,m a 400\,nm.

Un análisis cronológico de las marcas de tiempo del sistema de archivos
muestra que esta propuesta se escribió \emph{antes} que la de tablas
tabuladas del grupo 3 (15:11--15:28 frente a 15:45--16:32 del
2026-09-11), y que la cadena que efectivamente continuó hacia la revisión
final \codigo{-corregido} (con las correcciones de PMT) es la del grupo 3,
no esta. Como ambas modifican las mismas tablas de propiedades del mismo
material de forma mutuamente excluyente, \textbf{no es posible aplicar
ambas simultáneamente sin reintroducir una inconsistencia física}
(¿qué \codigo{RINDEX} usa el agua pura?). Se optó por integrar la rama
efectivamente continuada (grupo 3), dejando esta propuesta documentada
aquí como una alternativa evaluada y no adoptada en la versión final del
simulador. Su archivo fuente se conserva sin cambios en la carpeta de
insumos por si se desea retomar como línea de trabajo futura o como
comparación de sistemáticas.

\subsection{Grupo 5 --- Correcciones de la respuesta del PMT}

\begin{itemize}
  \item \textbf{\codigo{G4MPMTAction::ProcessHits}}: la respuesta del PMT
    ahora se activa únicamente cuando el fotón óptico cruza la frontera
    geométrica de entrada (\codigo{step->GetPreStepPoint()->GetStepStatus()
    == fGeomBoundary}); los fotones aceptados, rechazados por rango de
    energía y rechazados por eficiencia cuántica se terminan explícitamente
    (\codigo{track->SetTrackStatus(fStopAndKill)}), evitando que el mismo
    fotón genere más de un fotoelectrón en pasos repetidos dentro del
    volumen de Pyrex.
  \item \textbf{\codigo{OptDevice::GetQuantumEfficiency}}: la eficiencia
    cuántica nominal del PMT R5912 se interpola linealmente entre ocho
    puntos tabulados de 300 a 650\,nm, en vez de usar una función escalón
    de 50\,nm con saltos artificiales. La eficiencia de colección se aplica
    una única vez, sobre el resultado de la interpolación.
  \item \textbf{Eliminación de archivos de depuración por fotón}: se
    retiran las escrituras a \codigo{eficiencia-limpia*.txt} dentro de
    \codigo{IsPhotonDetected}/\codigo{GetQuantumEfficiency}, que generaban
    un archivo de una línea por cada fotón óptico simulado.
  \item \textbf{Material efectivo del PMT (\codigo{Pyrex})}: el rango de
    energía de su tabla óptica pasa de un intervalo fijo (2.00--4.00\,eV)
    a exactamente el mismo rango que \codigo{OptDevice} usa para la
    eficiencia cuántica (300--650\,nm, expresado en eV), y se documenta
    explícitamente que el fotocátodo es un absorbente efectivo: la
    eficiencia cuántica y de colección se muestrean una única vez en
    \codigo{G4MPMTAction}/\codigo{OptDevice}, no en la tabla de
    propiedades del material.
  \item \textbf{Geometría del WCD (\codigo{SaltyWCD::BuildDetector})}: se
    reescribe con capas explícitas y no superpuestas (agua/salmuera,
    Tyvek superior/inferior/lateral, acero superior/inferior/lateral) en
    vez del anidamiento de \emph{casings} del Meiga anterior, y se añaden
    verificaciones con \codigo{G4Exception} que abortan la construcción de
    la geometría si el medio activo no tiene \codigo{RINDEX}/\codigo{ABSLENGTH},
    si la superficie del Tyvek no tiene \codigo{RINDEX}/\codigo{REFLECTIVITY},
    o si el material efectivo del PMT no tiene \codigo{RINDEX}/\codigo{ABSLENGTH}.
\end{itemize}

\subsection{Reintegración del modelo térmico S($\alpha,\beta$)}
\label{sec:salphabeta}

Los cinco grupos de mejora se desarrollaron, según pudo determinarse
comparando \codigo{Materials-mejorado.cc} contra el Meiga anterior, a
partir de una copia de \codigo{Materials.cc} \emph{previa} a la
introducción del elemento \codigo{TS\_H\_of\_Water}: construían
\codigo{Water} con \codigo{elH} (hidrógeno libre) y densidad nominal de
Geant4 (1.0\,g/cm\textsuperscript{3}), en vez de con el elemento ligado y
la densidad medida a 20\,\textdegree C (0.99820\,g/cm\textsuperscript{3}).
Aplicar esa cadena de cambios tal cual habría eliminado silenciosamente el
tratamiento térmico S($\alpha,\beta$) ya validado en el Meiga anterior
--el objeto central de la variante \codigo{SalphaBeta-2026}-- para
\emph{todas} las soluciones de NaCl y para el agua pura.

En la integración final se restituyó \codigo{elTSHWater} en
\codigo{Materials::CreateElements}, se reconstruyó \codigo{Water} con el
mismo estado termodinámico explícito del Meiga anterior
(\codigo{kStateLiquid}, 293.15\,K, 1\,bar) usando \codigo{elTSHWater} en
vez de \codigo{elH}, y se propagó el mismo elemento ligado a los tres
materiales salinos nuevos del grupo 2
(\codigo{CreateSaltyWaterByFinalMassFraction} ahora recibe
\codigo{elTSHWater} como elemento de hidrógeno, y construye cada material
con el mismo estado termodinámico). De esta manera, \codigo{G4ThermalNeutrons}
--que se activa automáticamente sobre cualquier material que contenga un
elemento \codigo{TS\_*} reconocido, sin necesidad de cambios en
\codigo{G4MPhysicsList.cc}-- aplica la ley S($\alpha,\beta$) del agua de
forma consistente a los cuatro medios activos del WCD (agua pura y las
tres soluciones), tal como especifica el CHANGELOG original de la
variante SalphaBeta: \emph{“El H de las soluciones reutiliza la evaluación
h\_water; Na, Cl y O siguen el tratamiento estándar de G4NDL”}. Los
apéndices de la Sección~\ref{ap:diffs} muestran este parche como parte
del \emph{diff} de \codigo{Materials.cc}.

\section{Validación realizada y pendiente}
\label{sec:validacion}

\begin{itemize}
  \item \textbf{Realizada en este entorno} (sin instalación de Geant4
    disponible): verificación de balance de llaves y paréntesis en cada
    archivo modificado; búsqueda exhaustiva, en todo \codigo{src/}, de
    referencias a cada símbolo público añadido, renombrado o eliminado de
    \codigo{Materials}; confirmación de que \codigo{WCD.cc},
    \codigo{Hodoscope.cc}, \codigo{Mudulus.cc}, \codigo{Musaic.cc} y
    \codigo{Scintillator.cc} —que comparten la instancia \emph{singleton}
    de \codigo{Materials}— siguen compilando lógicamente contra la nueva
    interfaz; confirmación de que \codigo{G4WCDConstruction.cc} invoca
    \codigo{SaltyWCD::BuildDetector} con la misma firma.
  \item \textbf{Pendiente, a realizar en la máquina de simulación con
    Geant4 10.7.3 y G4NDL 4.6 instalados}: compilación completa
    (\codigo{cmake && make}), una corrida corta de verificación por cada
    una de las cuatro concentraciones de NaCl, e inspección de las nuevas
    tablas TSV generadas por \codigo{G4WCDSteppingAction} para confirmar
    que \codigo{capturas-neutrones.tsv} reporta $Z=1$ para el hidrógeno y
    que la fracción de capturas por debajo de 4\,eV es consistente con la
    obtenida en la validación original de la variante SalphaBeta.
\end{itemize}

\section{Archivos entregados}

Junto con este informe se entrega:

\begin{itemize}
  \item El árbol de código fuente completo de MEIGA con todas las mejoras
    integradas (\codigo{MEIGA-final/}), listo para compilar en la máquina
    de simulación.
  \item Copia de los ocho archivos originales del Meiga anterior que
    fueron modificados (carpeta \codigo{originales/}), y copia de su
    versión final integrada (carpeta \codigo{finales/}), para inspección
    directa lado a lado.
  \item El \emph{diff} unificado de cada uno de los ocho archivos
    (carpeta \codigo{diffs/} y Apéndice~\ref{ap:diffs} de este informe).
\end{itemize}

\section{Conclusiones}

Las cinco familias de mejoras revisadas en los informes técnicos previos
--registro paso a paso, densidades y composición de las soluciones de
NaCl, corrección óptica del Tyvek con tablas por concentración,
correcciones de la respuesta del PMT, y una propuesta óptica alternativa--
se integraron en una única base de código consistente derivada del Meiga
anterior. La integración requirió resolver tres inconsistencias no
triviales en los archivos de entrada (una rama óptica alternativa
incompatible, un archivo con nombre incorrecto, y un archivo huérfano de
nombre corrupto) y restituir el tratamiento térmico S($\alpha,\beta$) que
la cadena de archivos de mejora había perdido por partir de una copia
anterior a su introducción. El resultado conserva íntegramente el objetivo
original de la variante \codigo{SalphaBeta-2026} --dispersión térmica
realista del hidrógeno del agua y las soluciones de NaCl-- y añade sobre
él todas las correcciones ópticas y de respuesta del PMT revisadas,
dejando el simulador listo para compilar y validar en la máquina de
simulación.

\appendix

\section{Diffs unificados: Meiga anterior vs.\ versión integrada}
\label{ap:diffs}

Cada listado muestra el \emph{diff} unificado (\codigo{diff -u}) entre el
archivo original del Meiga anterior (\codigo{---}) y la versión final
integrada en este trabajo (\codigo{+++}). Las líneas que comienzan con
\codigo{-} se eliminaron; las que comienzan con \codigo{+} se añadieron.

\subsection{\texorpdfstring{\codigo{Applications/G4WCDSimulator/G4WCDSteppingAction.h}}{Applications/G4WCDSimulator/G4WCDSteppingAction.h}}
\lstinputlisting[style=diffstyle]{diffs/G4WCDSteppingAction.h.diff}

\subsection{\texorpdfstring{\codigo{Applications/G4WCDSimulator/G4WCDSteppingAction.cc}}{Applications/G4WCDSimulator/G4WCDSteppingAction.cc}}
\lstinputlisting[style=diffstyle]{diffs/G4WCDSteppingAction.cc.diff}

\subsection{\texorpdfstring{\codigo{G4Models/Materials.h}}{G4Models/Materials.h}}
\lstinputlisting[style=diffstyle]{diffs/Materials.h.diff}

\subsection{\texorpdfstring{\codigo{G4Models/Materials.cc}}{G4Models/Materials.cc}}
\lstinputlisting[style=diffstyle]{diffs/Materials.cc.diff}

\subsection{\texorpdfstring{\codigo{G4Models/SaltyWCD.cc}}{G4Models/SaltyWCD.cc}}
\lstinputlisting[style=diffstyle]{diffs/SaltyWCD.cc.diff}

\subsection{\texorpdfstring{\codigo{G4Models/G4MPMTAction.h}}{G4Models/G4MPMTAction.h}}
\lstinputlisting[style=diffstyle]{diffs/G4MPMTAction.h.diff}

\subsection{\texorpdfstring{\codigo{G4Models/G4MPMTAction.cc}}{G4Models/G4MPMTAction.cc}}
\lstinputlisting[style=diffstyle]{diffs/G4MPMTAction.cc.diff}

\subsection{\texorpdfstring{\codigo{Framework/Detector/OptDevice.cc}}{Framework/Detector/OptDevice.cc}}
\lstinputlisting[style=diffstyle]{diffs/OptDevice.cc.diff}

\section{Archivos de referencia sin cambios}

Los siguientes archivos, aportados como copias de referencia en la carpeta
\codigo{Correcciones-PMT/}, se compararon byte a byte contra el Meiga
anterior y resultaron \textbf{idénticos}, por lo que no requirieron ninguna
acción: \codigo{Detector.h}, \codigo{Detector.cc}, \codigo{OptDevice.h},
\codigo{OpticalPhysics.cc}, \codigo{OpticalPhysics.hh},
\codigo{G4MPMTAction.h} (copia sin sufijo) y \codigo{G4MPMTAction.cc}
(copia sin sufijo). También \codigo{G4Models/SaltyWCD.h} se mantuvo sin
cambios: ninguna de las revisiones de \codigo{SaltyWCD.cc} añadió
miembros nuevos a la clase, por lo que su interfaz pública
(\codigo{BuildDetector}) permanece igual.

\end{document}
